\chapter*{记号}
\markboth{记号}{记号}
\addcontentsline{toc}{chapter}{记号}
% Source PDF page 711 (printed p. 697)
\section*{一些特定的集合}
\begin{longtable}{@{}p{.28\textwidth}p{.66\textwidth}@{}}
$\R$ & 实数。\\
$\R^n$ & 实 $n$ 维向量（$n\times1$ 矩阵）。\\
$\R^{m\times n}$ & 实 $m\times n$ 矩阵。\\
$\R_+,\R_{++}$ & 非负实数、正实数。\\
$\C$ & 复数。\\
$\C^n$ & 复 $n$ 维向量。\\
$\C^{m\times n}$ & 复 $m\times n$ 矩阵。\\
$\Z$ & 整数。\\
$\Z_+$ & 非负整数。\\
$\mathbf S^n$ & 对称 $n\times n$ 矩阵。\\
$\mathbf S_+^n,\mathbf S_{++}^n$ & 对称半正定、对称正定 $n\times n$ 矩阵。
\end{longtable}

\section*{向量与矩阵}
\begin{longtable}{@{}p{.28\textwidth}p{.66\textwidth}@{}}
$\one$ & 所有分量均为一的向量。\\
$e_i$ & 第 $i$ 个标准基向量。\\
$I$ & 单位矩阵。\\
$X^T$ & 矩阵 $X$ 的转置。\\
$X^H$ & 矩阵 $X$ 的厄米转置（共轭转置）。\\
$\tr X$ & 矩阵 $X$ 的迹。\\
$\lambda_i(X)$ & 对称矩阵 $X$ 的第 $i$ 大特征值。\\
$\lambda_{\max}(X),\lambda_{\min}(X)$ & 对称矩阵 $X$ 的最大、最小特征值。\\
$\sigma_i(X)$ & 矩阵 $X$ 的第 $i$ 大奇异值。\\
$\sigma_{\max}(X),\sigma_{\min}(X)$ & 矩阵 $X$ 的最大、最小奇异值。\\
$X^\dagger$ & 矩阵 $X$ 的 Moore--Penrose 逆，又称伪逆。\\
$x\perp y$ & 向量 $x$ 与 $y$ 正交：$x^Ty=0$。\\
$V^\perp$ & 子空间 $V$ 的正交补。\\
$\diag(x)$ & 以 $x_1,\ldots,x_n$ 为对角元素的对角矩阵。\\
$\diag(X,Y,\ldots)$ & 以 $X,Y,\ldots$ 为对角块的分块对角矩阵。\\
$\rank A$ & 矩阵 $A$ 的秩。\\
$\mathcal R(A)$ & 矩阵 $A$ 的值域。\\
$\mathcal N(A)$ & 矩阵 $A$ 的零空间。
\end{longtable}

% Source PDF page 712 (printed p. 698)
\section*{范数与距离}
\begin{longtable}{@{}p{.28\textwidth}p{.66\textwidth}@{}}
$\norm{\cdot}$ & 一个范数。\\
$\norm{\cdot}_*$ & 范数 $\norm{\cdot}$ 的对偶范数。\\
$\norm{x}_2$ & 向量 $x$ 的欧几里得范数（$\ell_2$ 范数）。\\
$\norm{x}_1$ & 向量 $x$ 的 $\ell_1$ 范数。\\
$\norm{x}_\infty$ & 向量 $x$ 的 $\ell_\infty$ 范数。\\
$\norm{X}_2$ & 矩阵 $X$ 的谱范数（最大奇异值）。\\
$B(c,r)$ & 中心为 $c$、半径为 $r$ 的球。\\
$\dist(A,B)$ & 集合（或点）$A$ 与 $B$ 之间的距离。
\end{longtable}

\section*{广义不等式}
\begin{longtable}{@{}p{.28\textwidth}p{.66\textwidth}@{}}
$x\preceq y$ & 向量 $x$ 与 $y$ 之间的逐分量不等式。\\
$x\prec y$ & 向量 $x$ 与 $y$ 之间的严格逐分量不等式。\\
$X\preceq Y$ & 对称矩阵 $X$ 与 $Y$ 之间的矩阵不等式。\\
$X\prec Y$ & 对称矩阵 $X$ 与 $Y$ 之间的严格矩阵不等式。\\
$x\preceq_K y$ & 正常锥 $K$ 诱导的广义不等式。\\
$x\prec_K y$ & 正常锥 $K$ 诱导的严格广义不等式。\\
$x\preceq_{K^*}y$ & 对偶广义不等式。\\
$x\prec_{K^*}y$ & 对偶严格广义不等式。
\end{longtable}

\section*{拓扑与凸分析}
\begin{longtable}{@{}p{.28\textwidth}p{.66\textwidth}@{}}
$\operatorname{card}C$ & 集合 $C$ 的基数。\\
$\operatorname{int}C$ & 集合 $C$ 的内部。\\
$\relint C$ & 集合 $C$ 的相对内部。\\
$\cl C$ & 集合 $C$ 的闭包。\\
$\operatorname{bd}C$ & 集合 $C$ 的边界：$\operatorname{bd}C=\cl C\setminus\operatorname{int}C$。\\
$\conv C$ & 集合 $C$ 的凸包。\\
$\aff C$ & 集合 $C$ 的仿射包。\\
$K^*$ & 与 $K$ 对应的对偶锥。\\
$I_C$ & 集合 $C$ 的示性函数。\\
$S_C$ & 集合 $C$ 的支撑函数。\\
$f^*$ & 函数 $f$ 的共轭函数。
\end{longtable}

\section*{概率}
\begin{longtable}{@{}p{.28\textwidth}p{.66\textwidth}@{}}
$\E X$ & 随机向量 $X$ 的期望。\\
$\operatorname{prob}S$ & 事件 $S$ 的概率。\\
$\operatorname{var}X$ & 标量随机变量 $X$ 的方差。\\
$N(c,\Sigma)$ & 均值为 $c$、协方差（矩阵）为 $\Sigma$ 的高斯分布。\\
$\Phi$ & $N(0,1)$ 随机变量的累积分布函数。
\end{longtable}

% Source PDF page 713 (printed p. 699)
\section*{函数与导数}
\begin{longtable}{@{}p{.28\textwidth}p{.66\textwidth}@{}}
$f:A\to B$ & $f$ 是定义在集合 $\dom f\subseteq A$ 上、取值于集合 $B$ 的函数。\\
$\dom f$ & 函数 $f$ 的定义域。\\
$\epi f$ & 函数 $f$ 的上图。\\
$\nabla f$ & 函数 $f$ 的梯度。\\
$\nabla^2 f$ & 函数 $f$ 的 Hessian 矩阵。\\
$Df$ & 函数 $f$ 的导数（Jacobi）矩阵。
\end{longtable}
% Source PDF page 714 (printed p. 700) is blank.
